\documentclass[a4paper]{report}
\usepackage[utf8]{inputenc}
\usepackage[T1]{fontenc}
\usepackage[french]{babel}
\usepackage{setspace}
\usepackage{graphicx}
\usepackage{wrapfig}
\usepackage{lscape}
\usepackage{bookmark}
\usepackage{rotating}
\usepackage{epstopdf}
\usepackage{subcaption}
\usepackage{float}
\usepackage{eurosym}
\usepackage[headheight=13pt,top=3cm, bottom=2cm, left=1cm, right=1cm]{geometry}
\usepackage{pdfpages}
\usepackage{fancyhdr}
\pagestyle{fancy}
\usepackage{hyperref}
\usepackage{array,multirow,makecell}
\setcellgapes{3pt}
\setlength{\parindent}{1.1cm}
\makegapedcells
\newcolumntype{L}[1]{>{\raggedright\arraybackslash }b{#1}}
\newcolumntype{R}[1]{>{\raggedleft\arraybackslash }b{#1}}
\newcolumntype{C}[1]{>{\centering\arraybackslash }b{#1}}
\renewcommand{\footrulewidth}{1pt}
\fancyfoot[L]{ENSIAS}
\fancyfoot[C]{\textbf{\thepage}}
\fancyfoot[R]{Année universitaire: 2020/2021}
\begin{document}
\begin{titlepage}
    \begin{center}
        \begin{figure}[!h]
            \vspace{- 2 cm}
            \hspace{ 0 cm}
            \includegraphics[width=9em]{images/ensias.jpeg}
        \end{figure}
        \begin{figure}[!h]
            \vspace{- 3.34cm}
            \hspace{14 cm}
            \includegraphics[width=10em]{images/um5.jpeg}
        \end{figure}
    \end{center}

    \begin{center}
        \noindent \hspace{ 0.3 cm }\Huge \textbf{Rapport de Stage Fin de 2ème Année}
        \vspace*{0.1cm}

        \vspace*{0.1cm}
        \begin{center}
            \rule{0.9\linewidth}{1pt}
        \end{center}
        \begin{center}
            \noindent \hspace{ 0.4 cm}{\large \textsc{ Création d'une platform Web pour la gestion et le suivi des contrats du transport délégués}}\\
        \end{center}


        \vspace*{0.5cm} \noindent \hspace{ -0.5 cm} \large
        \begin{figure}[H]
            \begin{center}
                \includegraphics[width=3.5in,height=3in]{images/logo-uir2.jpg}
            \end{center}
        \end{figure}
        \raggedright
        {\textbf{\emph{Préparé par: KOTBI Abderrahamane}}}
        \raggedleft\rule{0.7\linewidth}{2pt}\\
        \raggedleft{\textbf{\emph{Membre Jury: *****}}}\\
        \Large\emph{Année universitaire : 2020-2021}

    \end{center}
\end{titlepage}
\pagenumbering{arabic} \setcounter{page}{1}
\addcontentsline{toc}{chapter}{Remerciements}
\begin{doublespace}
    \begin{center}
        \vspace*{1cm}

        \textbf{\huge{Remerciements}}

    \end{center}
\end{doublespace}
\begin{doublespace}

    \textit{
        Avant d'aborder la description des parties importantes du projet, j'aimerai tout d’abord exprimer ma gratitude
        intense à toute personne qui a contribué énormément dans l'élaboration et la réalisation de ce travail. Je commence
        ainsi par offrir mes remerciements à l'intégralité des personnes travaillants au sein de l’Ecole Nationale supérieure
        d’informatique et d’analyse des systèmes. Ces bonnes hommes et femmes qui ne cessent de j'aidere de j'prépae, moi et
        mes collègues, pour de telles expériences afin de bien s'integrer dans le cadre réelle du marché de travail.
        Toute ma gratitude et profonde reconnaissance s’adressent à à tous mes encadrants de stages: Mme.\textbf{***f***} ****,
        M..\textbf{***m***} ****   . Je vous remercie énormément pour m'avoir garanti un encadrement de qualité pour bien mener
        et assurer la réalisation de ce travail dans les meilleures des conditions.
    }

\end{doublespace}

\clearpage
\addcontentsline{toc}{chapter}{Résumé}
\begin{doublespace}
    \begin{center}
        \vspace*{1cm}

        \textbf{\huge{Résumé}}

    \end{center}
\end{doublespace}
\begin{doublespace}

    Ce document représente une étude synthétique du travail réalisé dans le cadre de mon stage d'été de deuxième année. J'ai effectué mon stage au sein de
    l'université internationale de Rabat \textbf{UIR}. Ce dernier a duré les deux mois juillet et août de l'année 2021. L’objectif principal de ce travail
    consiste à développer une platform Web pour la gestion et le suivi des contrats du transport délégués. D'ailleurs, le travail effectué dans le cadre
    de ce projet m'a représenté la bonne chance pour appliquer mes savoirs et savoirs-faire acquis durant cette année. En ce sens, la réalisation du projet
    a passé en trois étapes principal: l'élaboration d'une analyse détaillée des besoins du client, conception des différentes composantes de l'application,
    et enfin l'implémentation des objectifs soulignés. Par conséquent, j'ai eu la chance d'utiliser plusieurs téchnologies, outils et concepts
    qui m'ont aidé à développer cette solution.

\end{doublespace}
\begin{doublespace}

    Le reste de ce rapport est organisé comme suit : le chapitre I est dédié à la présentation de
    l’organisme d’accueil et du client. Quant au chapitre II, il traite l’analyse du probléme ainsi que
    la spésification des besoins. Ensuite le chapitre III présente l´étude conceptuelle. Enfin le chapitre
    IV traite la partie réalisation.

    \textbf	{\\Mots clés :\\ Gestion, Suivi, Donnée, Projet, Contrats délégués, Transport urbain, Centralisation, Application low-code.}

\end{doublespace}
\clearpage
\addcontentsline{toc}{chapter}{Abstract}
\begin{doublespace}
    \begin{center}
        \vspace*{1cm}
        \textbf{\huge{Abstract}}
    \end{center}
\end{doublespace}
\begin{doublespace}

    This document represents a synthetic study of the work carried out as part of my second-year summer internship.
    I did my internship at the international university of Rabat \textbf{UIR}. The latter lasted the two months
    July and August of the year 2021. It aims to develop a web platform for the management and monitoring of
    delegated transport contracts. Moreover, the work carried out in this project represented good oportunity
    for me to apply and sharpen my skills and my knowledge acquired during this year. Thus, the realization
    of the project has gone through three main stages: the extraction of a detailed analysis of the client's
    needs, the design of the application's different components, and finally the implementation of the outlined objectives.
    Therefore, I had the chance to use several technologies, tools, and concepts that helped me develop this solution.

    \textbf{\\Keywords:\\Management, Monitoring, Data, Project, Delegated contracts, Urban transport, Centralization, Low-code application.}

\end{doublespace}

\newpage

\renewcommand{\contentsname}{Table de matières}
\setcounter{tocdepth}{4}
\tableofcontents

\cleardoublepage

\addcontentsline{toc}{chapter}{\listfigurename}
\listoffigures
\newpage
\addcontentsline{toc}{chapter}{Introduction générale}
\begin{doublespace}
    \begin{center}
        \vspace*{1cm}
        \textbf{\huge{Introduction générale}}
    \end{center}
\end{doublespace}
\begin{doublespace}
    \renewcommand{\headrulewidth}{1pt}
    \fancyhead[R]{\textbf{Introduction générale}}
    \fancyhead[L]{\hspace*{5cm}}
    Dans le cadre de la modernisation de la gestion des secteurs publics collectifs
    marocains, l’état a créé des contrats de gestion déléguées. Parmi ces dernières, les
    contrats de gestion de transport urbain. Placer le citoyen au cœur des politiques et
    services publics en constitue un objectif fondamentale. En conséquence, le contrat
    vise à assurer des services publics de qualité, accessibles à tous les
    citoyens sans discrimination, à des coûts maitrisés tenant compte du pouvoir
    d’achat. D'où  l'importance de la bonne gestion et le suiv de ces contrats, pour garder
    les droits et les attentes de la population.
    Par ailleurs, l’état a pris ce sens pour réaliser ses objectifs d’ouverture sur l’économie
    mondiale, le soutien des entreprises publiques, et l’amélioration de ses services.
    En parallèle avec ceci, l’état a essayé de gagner les défis de la régionalisation.
    Par conséquent, la gestion et le suivis des contrats de transport a devenu une tâche
    fondamentale. Or, effectuer cette tâche en utilisant la méthode classique, par plusieurs documents,
    ne représente plus une bonne solution, vue le nombre de lecture et de vérification
    des données importantes dans les documents nombreux. En plus, cela ne permet pas d’en
    tirer profits pour prendre les bons décisions facilement. Ainsi, je propose une
    solution simple pour la gestion et la représentation des données en question.
\end{doublespace}

\newpage

\chapter{Cadre générale du projet}

\renewcommand{\headrulewidth}{1pt}
\fancyhead[R]{\textbf{Chapitre \thechapter: Cadre générale du projet}}
\fancyhead[L]{\hspace*{5cm}}
\begin{doublespace}
    Dans ce chapitre j'entame le projet dans son cadre général : présentation de l’organisme d’accueil \textbf{UIR},
    et de client, présentation de l'idée de l'application, présentation du contexte du projet, de la problématique et des objectifs fixés,
    et de la démarche de la réalisation que j'ai adopté.
\end{doublespace}
\begin{doublespace}
    \section{Présentation de l’organisme d’accueil}
    \subsection{UIR}
    L’Université Internationale de Rabat ou \textbf{UIR} est une université privée fondée en 2010 sous contrat avec l’État marocain.
    L'\textbf{UIR} concrétise ainsi le premier partenariat public-privé dans le domaine de l'enseignement supérieur au Maroc.
    Poursuivant l'objectif d'accompagner le Royaume du Maroc dans son développement, l'\textbf{UIR} a développé un catalogue
    de formation de haut niveau en adéquation avec les différentes stratégies impulsées par le Maroc (Plan solaire marocain
    , plan d'accélération industrielle, plan de digitalisation, etc.). Sûrtout sur le plan informatique, l'\textbf{UIR} demeure
    un partenaire important de l'état. Elle a plusieurs contributions en terme d'éducation et de préparation des cadres, ainsi
    qu'en terme de consulting, de recherche, de résolutions des problèmes émergeants, et en terme d'inovation.
    \begin{figure}[H]
        \begin{center}
            \includegraphics[scale=0.1]{images/logo-uir.jpg}
            \caption{Organisme d'acceuil}
        \end{center}
    \end{figure}
    \subsection{Partenaire de consulting technique: Dyn IT}
    \textbf{DYN IT MAROC} est le résultat de l’expérience de plusieurs Consultants en IT et Education.
    Nous conseillons et proposons des services en Ingénierie Informatique. Notre mission est de
    vous faire collaborer efficacement. Elle s'engage à aider les entreprises à mieux performer
    en fournissant des solutions de collaboration flexibles, évolutives et surtout abordables.
    Elle offrent des services de consultation et de formation en téchnologies Microsoft. En fait,
    \textbf{DYN IT} est un partenaire Microsoft.
    \begin{figure}[H]
        \begin{center}
            \includegraphics[scale=0.6]{images/dynit.png}
            \caption{Organisme d'acceuil}
        \end{center}
    \end{figure}
    \section{Présentation du Client}
    La Direction Générale des Collectivités Territoriales \textbf{DGCT} est chargée de la préparation des décisions du ministre de l’Intérieur,
    dans le cadre des attributions qui lui sont conférées en vertu des textes législatifs et réglementaires relatifs aux collectivités
    territoriales, et du suivi de leur exécution. Elle assure également l’appui et l'accompagnement juridique, technique et financier
    des collectivités territoriales, des instances qui en relèvent, des établissements de coopération intercommunale et des groupements
    des collectivités territoriales. Elle est chargée également, en coordination avec les départements et organismes concernés, de
    concourir au développement territorial. Bref, ses missions sont en gros:
    \begin{itemize}
        \item Planification et développement territorial.
        \item Assistance des réseaux publics locaux, et des institutions locales.
        \item Suivi juridique et gestion des services locaux.
        \item Amélioration de la mobilité urbaine et du transport.
        \item Développement des compétences et transformation digitale.
        \item Accompagnement financier des collectivités territoriales.
        \item Coopération décentralisée.
    \end{itemize}

    \begin{figure}[H]
        \begin{center}
            \includegraphics[scale=0.27]{images/logo-fr.png}
            \caption{Client}
        \end{center}
    \end{figure}
    \newpage
    \section{Problématique}
    Dans le cadre de la bon gestion de territoire et des aménagements publics, plusieurs
    contrats de gestion déléguée du transport urbain sont signés. Ces derniers comportent
    plusieurs actes qui forment un ensemble de règles à respecter entre un délégant qui
    sont une ou plusieurs communes et un délégataire qui est une société. En addition,
    des pénalités sont imposeées en cas de non respect des règles posées dans le contrat.
    Par conséquent, le suivi de la réalisation et du respect de ses actes est un devoir
    juridique. D’où l'importance d'avoir une solution pertinente pour veiller sur le
    respect des conditions prédéfinis.
    \\\textbf{Alors, comment peut on améliore les opérations de suivis et de traitements des données en questions?}
    \section{Solution et Objectifs du projet}
    La première étape de la conduite de projet est sans doute la plus évidente, mais peut-être
    aussi la plus cruciale. En effet, sans objectifs bien définis, il est difficile de savoir où votre projet
    va vous mener. Pour cela, avant de commencer dans la conception et la réalisation de ce projet,
    il faut tout d’abord fixer l’objectif principal de manière à développer la ligne d’actions à mener.
    Je cherche à concevoir et construire une application Web qui permet de saisir, d'enregistrer,
    et de traiter les données qui concernent les contrats de gestion de transport. La
    solution proposée est composée de deux partie: une première partie qui concerne la
    saisi et l'enregistrement des données sous la forme d'une plateforme web liée à une
    base de donné, et une deuxième partie qui concerne l'analyse des données.
    Bref, je cherche une solution pour:
    \begin{itemize}
        \item Centraliser l’information contenue dans les contrats de gestion déléguée et
              leurs avenants le cas échéant.
        \item Centraliser les documents exigés par le contrat que l’opérateur est tenu de
              fournir périodiquement à l’autorité délégante.
        \item Permettre aux utilisateurs de partager en interne des informations et des
              documents.
        \item Assurer le suivi de la gestion des contrats de transport urbain.
    \end{itemize}
    \newpage
    \section{Récapitulatif }
    Dans ce chapitre introductif, j'ai pu décrire le conteste général du projet, et
    déterminer son objectif principal. En addition, les aspects que je vai
    avoir besoin lors de la réalisation de cette application.
    Le chapitre qui suit consiste la phase d'analyse des besoins du projet.
    \newpage
    \chapter{Analyse et conception}

    \renewcommand{\headrulewidth}{1pt}
    \fancyhead[R]{\textbf{Chapitre \thechapter: Analyse et conception}}
    \fancyhead[L]{\hspace*{5cm}}
    \begin{doublespace}

        Ce chapitre représente le point de départ de mon travail. Premièrement, j'analyse et
        je spécifie les besoins du projet. Ensuite, j'identifie les différents acteurs. Et enfin, je
        modélise le tout dans un diagramme des cas d’utilisation général qui sera notre file conducteur
        durant la prochaine phase.
        \section{Analyse}
        \subsection{Introduction à l'analyse}
        Certes, la gestion et le suivis des données en utilisant des solutions Web ne date pas d'hière. Mais,
        elle demeure une solution optimisante et interéssante en terme de temps et de productivité. Sûrtout,
        dans un tel cas où les informations sont, à la fois, cruciale et nombrueses. Ainsi, la solution
        proposée premièrement par la \textbf{DGST} est d'avoir une plateforme Web pour faciliter la collècte,
        le traitement, la présentation, et la centralisation de l'information. Sur la même longueur d'onde, j'ai
        travaillé sur la solution proposée avec une équipe de la \textbf{DGST} pour apporter les bonnes
        fonctionnalités répondantes aux besoins de cette dernière.
        \subsection{Étude de l’existant}
        Après une étude du marché, on a trouvé une diversité des sites web et des applications dédiées à la
        gestion et au suivi.
        Comme un exemple, considérer les suivants deux applications:\\
        \textbf{\large Jumia:}\\
        Jumia est une entreprise de commerce en ligne active sur le marché africain qui a été fondée en 2012. Jumia, la plateforme en ligne marché qui relie les vendeurs et acheteurs en fournissant une solution logistique qui comprend l'entreposage, la distribution, et la livraison de colis, de même que d'un  service paiement responsable.\\
        \textbf{Avantages:}
        \begin{itemize}
            \item Une bonne gestion des promotions.
            \item Un service de paiement en ligne ou à la livraison.
            \item Une organisation et classification des produits.
        \end{itemize}
        \textbf{Inconvénients:}
        \begin{itemize}
            \item Délais de livraison très importants.
            \item Existence des redondance (Produits identiques avec deux a trois prix différents).
        \end{itemize}
        \textbf{\large AliExpress:}\\
        \begin{doublespace}
            Aliexpress.com est un site de commerce en ligne du Groupe Alibaba spécialisé dans la vente
            de produits à prix bas, aux particuliers (BC) et à l’international.\\
            \textbf{Avantages:}
            \begin{itemize}
                \item Une livraison à l’échelle internationale.
                \item Une protection pour l’acheteur et des paiements sécurisés.
                \item Une possibilité de voir les retours clients.
            \end{itemize}
            \textbf{Inconvénients:}
            \begin{itemize}
                \item Une durée de livraison élevée.
                \item Peu de controle sur la marchandise.
            \end{itemize}
            \newpage
            \section{Spécification des besoins:}
            \begin{doublespace}
                \subsection{Spécification des acteurs:}
                \begin{itemize}
                    \item L’administrateur: c'est la personne chargée d'affecter les différents rôles et de
                          gérer les comptes et les produits ainsi que les réclamations.
                    \item  Le client: Il consulte le catalogue et les offres  spéciaux.Il est possible pour lui de mener des recherches.En même que la capacité de voir les promotions.
                    \item Le visiteur: Un utilisateur ordinaire qui n'est pas censé à d'avoir un compte ou  de s'authentifier, mais il peut voir les produits et les nouveautés.
                \end{itemize}
                \subsection{Besoins fonctionnels:}
                \begin{itemize}
                    \item Gestion des utilisateurs: L'administrateur doit avoir une section d'où il pourra faire une gestion des clients, ainsi de leurs détails et activités.
                    \item Gestion des produits: Les produits mis sur la plateforme de manière bien ordonnée.
                    \item Gestion des transaction : l’administrateur peut consulter les transaction dans l'application.
                    \item Gestion des ventes: Outil de prévision qui doit fournir des graphes fiables.
                \end{itemize}
                \subsection{Les besoins non fonctionnels:}
                Les besoins fonctionels sont basiques pour un fonctionnement correcte et une réponse fiable aux besoins des utilisateurs, mais il y a des autres besoins qui tendent à améliorer la performance et la qualité de l'application pour une utilisation plus adéquate.
                \begin{itemize}
                    \item Fiabilité de la plateforme: L’application doit fonctionner sans erreur.
                    \item Ergonomie, souplesse et confort d’utilisation: Pour faciliter l’utilisation, notre plateforme doit offrir une interface unifiée, conviviale et érgonomique.
                    \item Gain de temps: Notre application doit optimiser les traitements pour avoir un temps de réponse minimale.
                    \item Maintenabilité et sociabilité: Notre code source doit être lisible, commenté,compréhensible afin d’assurer son état évolutif et extensible par rapport aux besoins des utilisateurs.
                    \item Sécurité: Notre plateforme doit  être très authentique en ce qui concerne les informations personnelles des utilisateurs.
                \end{itemize}
                \section{Conception:}
                \begin{doublespace}
                    Après la prise en charge de la question, j'avonseordé le principe de la création d'un site e-commerce qui répond aux normes et aux modalités existantes dans les applications du commerce digitales, avec l'daptation des nouvelles technique qui touchent la globalité de la solution, avec la possibilité d'une connexion à chaque tentative d'accès.

                    \subsection{Cas d'utilisation globale:}
                    \begin{figure}[H]
                        \begin{center}
                            %  \includegraphics[scale=0.7]{images/pfa2 utilisation globale.png}
                            \caption{Cas d'utilisation globale}
                        \end{center}
                    \end{figure}
                    \newpage
                    Dans cette figure on fait la segmentation des besoins principaux dont l'application doit répondre, on distingue les trois types des utilisateurs Administrateur, client et visiteur: \\Les deux premiers types sont les seuls à avoir des comptes avec authentification de soi, la fonctionnalité de l'administateur se focalise sur tout ce qui est comparable à la gestion.Toutefois, le client peut effectuer des achats ainsi qu'il peut évaluer son expérience et ses produits. \\ L'application est aussi ouverte pour les visiteurs ordinaux, ils peuvent avoir un panier, consulter les produits à la phase de la création des comptes ils peuvent devenir des clients.

                    \subsection{Cas d'utilisation de l'administrateur:}
                    \begin{figure}[H]
                        \begin{center}
                            %  \includegraphics[scale=0.7]{images/admin cas.png}
                            \caption{Cas d'utilisation de l'administrateur}
                        \end{center}
                    \end{figure}
                    Dans cette figure on illustre le cas de l'utilisation de l'administrateur de façon plus détaillée, comme c'est déjà mentionné l'administrateur après son authentification, il prend la responsabilité de la gestion des clients, leurs transactions, ainsi qu'aux produits du magazins.Il fait la mise à jour, la suppression ou l'ajout des produits pour une gestion des produits plus fluide.\\ N'oubliant pas qu'il aura la possibilité de la consultation des informations des comptes clients leurs détails et leurs transactions effectuées.

                    \subsection{Cas d'utilisation du client:}
                    \begin{figure}[H]
                        \begin{center}
                            %  \includegraphics[scale=0.7]{images/client cas (1).png}
                            \caption{Cas d'utilisation du client}
                        \end{center}
                    \end{figure}

                    L'utilisateur client est est l'acteur principale qu'on vise lors de la conception de l'application, il a un compte personnel d'où il peut effectuer plusieurs fonctionnalités dès qu'il s'authentifie il n'est plus visiteur. Lors de la consultation de son compte il peut faire des modifications sur le plan de ses données, pour les achats cette fonctionnalité est triviale ainsi disponible selon les critères de chaque produit.\\D'ou il peut faire l'évalution et la consultation de chaque produit dans sa commande. \\ La communication au client et le service de la vente forment des traits principaux lors de l'achat et la vente, ici le client peut toujours rapporter ses interrogations et les abus d'utilisation à l'administrateur.
                    \newpage
                    \subsection{Cas d'utilisation visiteur:}
                    \begin{figure}[H]
                        \begin{center}
                            %  \includegraphics[scale=0.7]{images/client cas (2).png}
                            \caption{Cas d'utilisation du visiteur}
                        \end{center}
                    \end{figure}

                    Le client avant qu'il obtient un compte à lui, il reste un visiteur ordinaire, avec diffèrents aspects d'utilisation: chaque visiteur aura un panier dont il pourra effectuer ses achats, consulter les produits qui existent dans son panier ainsi qu'il peut ajouter des produits dans ce panier.\\
                    Il peut à tout moment contacter l'administrateur pour les renseignements, les modalités de paiement des commandes et de livraison colis.

                    \subsection{Diagramme de séquence Authentfication Administrateur/ client:}
                    \begin{figure}[H]
                        \begin{center}
                            %  \includegraphics[scale=0.6]{images/sec auth (2).png}
                            \caption{Diagramme de séquence Authentfication Administrateur/ client}
                        \end{center}
                    \end{figure}
                    Selon l'acteur, le système approuve les sessions d'utilisation des deux acteurs Administrateur ou client: Lors de l'authentification le système exige la validité des identifiants avec la vérification de la base de données, on approuve la connexion à la base de ces résultats.
                    \subsection{Diagramme de séquence création d'un compte client:}
                    \begin{figure}[H]
                        \begin{center}
                            %  \includegraphics[scale=0.6]{images/sec auth (3) (1).png}
                            \caption{Diagramme de séquence création d'un compte client}
                        \end{center}
                    \end{figure}
                    La séquence  qui précède décrit la création des comptes effectuée par les visiteurs pour devenir des clients. Le visiteur entre ses données dans la case convenable, si ces données sont qualifiées valides alors le compte est valable. Le visiteur est un client à ce moment.
                    \subsection{Diagramme de séquence ajout au panier:}
                    \begin{figure}[H]
                        \begin{center}
                            %  \includegraphics[scale=0.6]{images/sec auth (4).png}
                            \caption{Diagramme de séquence ajout au panier}
                        \end{center}
                    \end{figure}
                    Pendant chaque achat, les clients consultent les produits, font leurs recherches puis arrivent à choisir les produits qui leurs convient, d'où l'idée d'avoir un panier comme l'expérience au magazin réellement. \\Le client ajoute des produits à son panier, mais pour qu'il puisse accomplir son achat, l'application nécessite une phase d'enregistrement des données clients, la vérification des données si favorable permet au client de compléter son achat et d'accéder au paiement finalement.

                    \subsection{Diagramme de séquence Ajouter ou Supprimer ou modifier les catégories:}
                    \begin{figure}[H]
                        \begin{center}
                            %  \includegraphics[scale=0.6]{images/sec auth (5).png}
                            \caption{Diagramme de séquence modifications des catégories}
                        \end{center}
                    \end{figure}
                    Pour assurer la bonne gestion des diffèrents produits sur notre plateforme, on a fait de sorte de classifier les produits sous des catégories, chacune contient un genre différent des produits:\\ L'administrateur est l'acteur en charge de faire des telles modifications comme la suppression, l'ajout ou la mise à jour des catégories dont on inscrit les produits. L'administrateur s'authentifie au système, avec un retour de validation il effectue la modification, suppression ou ajout. Cette action est envoyée à la base de donnée pour un enregistrement des changements, la base de données traite ces actions le retour donne l'application de ces changements sur plateforme réelle.
                    \section{Structure conceptuelle de l'application:}
                    \begin{figure}[H]
                        \begin{center}
                            %  \includegraphics[scale=0.6]{images/Capp.PNG}
                            \caption{Structure des catégories}
                        \end{center}
                    \end{figure}
                    Pour bien expliquer la vision conceptuelle de l'application réelle , on consacré une phase de l'étude pour faire décider la structure de l'interface utilisateur  qui nécessite la conception des  sections suivants  disponible lors de la phase d'acceuil.
                    \begin{itemize}
                        \item \textbf{Ordinateurs portables:} Où les acteurs trouvent la condition exacte du produit (Nouveau ou non ).
                        \item \textbf{Accessoires:} L'utilisateur peut savoir le niveau du stockage de chaque outil ainsi que l'ergonomie de ces accessoires les souris, claviers ainsi les cables.
                        \item \textbf{Contact/Plus d'informations:} A partir de celles-ci les acteurs en question peuvent contacter l'administrateur.
                        \item \textbf{Login/Register:} C'est la phase de l'authentification aux comptes.
                        \item \textbf{Panier:} C'est l'entité destinée aux clients lors de l'achat des produits.

                    \end{itemize}


                    \newpage
                    \section{Récapitulatif:}
                    Face aux conditions actuelles de la distanciation sociale et la fermeture des magazins et des espaces makro, cela a bien incarné l'idée de l'achat en ligne d'où on a eu l'idée de la création d’une application web de l'e-commerce spécialisée en vente des gadgets électronique. Cette partie aborde encore toutes les besoins des utilisateurs dont le site doit répondre : celles de l’administrateur, le client  et même du visiteur.

                    Dans la première section on a essayé de bien analyser à fond les besoins primordiales des utilisateurs qu'on a spécifié lors de la phase de l'argumentation des besoins, où on a cité trois types d'acteurs avec les typologies de leurs besoins, sous forme plus graphique on a utilisé le langage UML pour décrire tout les cas d'utilisation possible à chaque utilisateur de façon globale puis  détaillée pour bien déterminer les limites d'action de chaque utilisateur.
                    Ensuite on a décider de rendre les actions plus dynamiques, avec des diagrammes de séquence, qui font la description des interactions entre le système et les acteurs lors de chaque cas d'utilisation.





                    \chapter{Réalisation:}
                    \fancyhead[R]{\textbf{Chapitre \thechapter:Réalisation}}
                    \fancyhead[L]{\hspace*{5cm}}
                    \begin{doublespace}

                        \section{Technologies de travail:}
                        Ce chapitre citera toutes les technologies et les outils utilisés pour la mise en œuvre de l'application Laptop+ .
                        \subsection{Langage Python:}
                        \begin{figure}[H]
                            \raggedleft{
                                %  \includegraphics[scale=0.3]{images/pythob.png}
                            }
                        \end{figure}
                        Python est un langage de programmation largement utilisé, interprété, orienté objet et de haut niveau avec une sémantique dynamique, utilisé pour la programmation à usage général. Il a été créé par Guido van Rossum et sorti pour la première fois le 20 février 1991. C'est un langage simple et intuitif tout aussi puissant que ceux des grands concurrents,open source , pour que chacun puisse contribuer à son développement , il est aussi compréhensible
                        adapté aux tâches quotidiennes , permettant des temps de développement courts.

                        \newpage

                        \subsection{Frame work Django:}
                        \begin{figure}[H]
                            \raggedleft{
                                %  \includegraphics[scale=0.3]{images/django.png}
                            }
                        \end{figure}
                        Django est un framework Web Python de haut niveau qui encourage un développement rapide et une conception propre et pragmatique. Construit par des développeurs expérimentés, il prend en charge une grande partie des tracas du développement Web, vous pouvez donc vous concentrer sur l'écriture de votre application sans avoir à réinventer la roue. C'est gratuit et open source.

                        \subsection{L'éditeur de texte Atome:}
                        \begin{figure}[H]
                            \raggedleft{
                                %  \includegraphics[scale=0.3]{images/atome.png}
                            }
                        \end{figure}
                        Atom est un éditeur de texte et de code source gratuit et open-source pour macOS , Linux et Microsoft Windows avec prise en charge des plug-ins écrits en JavaScript et Git Control intégré , développé par GitHub . Atom est une application de bureau construite à l'aide des technologies Web .La plupart des packages d'extension ont des licences de logiciel libre et sont construits et maintenus par la communauté.
                        \newpage
                        \subsection{ Les graphiques Chart.Js:}
                        \begin{figure}[H]
                            \raggedleft{
                                %  \includegraphics[scale=0.4]{images/cahrtjs.png}
                            }
                        \end{figure}
                        Chart.js est une bibliothèque JavaScript open source gratuite pour la visualisation de données , qui prend en charge 8 types de graphiques : bar , line , area , pie ( donut ), bubble , radar , polar et scatter. Créé par le développeur Web  Nick Downie, il est maintenant maintenu par la communauté et est la deuxième bibliothèque de graphiques JS la plus populaire sur GitHub par le nombre d'étoiles après D3.js , considéré comme beaucoup plus facile à utiliser mais moinspersonnalisable que ce dernier. Chart.js est rendu dans un canevas HTML5 et est largement considéré comme l'une des meilleures bibliothèques de visualisation de données.
                        \subsection{Databases Sqlite:}
                        \begin{figure}[H]
                            \raggedleft{
                                %  \includegraphics[scale=0.4]{images/sqlite.jpg}
                            }
                        \end{figure}
                        SQLite est une bibliothèque in-process qui implémente un moteur de base de données SQL transactionnel autonome, sans serveur, à configuration zéro . Le code de SQLite est dans le domaine public et est donc libre d'utilisation à toutes fins, commerciales ou privées.
                        \newpage
                        \section{L'application Web Laptp+ :}
                        L'application Laptop+ est une application web e-commerce spécialisée dans la vente des ordinateurs portables et les accessoires en relation avec ces derniers. \\Laptop+ expose à la vue une plateforme d'achat est si innovante et puissante, son design relaxant et l’explicité de son contenu crée un espace compatible à tous les utilisateurs.

                        \section{ Visualisation de la base de données:}
                        \subsection{ Table des catégories: }
                        \begin{figure}[H]
                            %  \includegraphics[scale=0.5]{images/WhatsApp Image 2021-06-21 at 15.34.09.jpeg}
                            \caption{Table des catégories}
                        \end{figure}
                        Cette figure donne l'aspect technique des catégories citées à l'étude
                        fonctionnelle , chaque produit a son propre Id, le lien exact c'est à dire le slag ainsi qu'il appartient à une catégorie de produits spécifique.
                        \subsection{ Table des produits: }
                        \begin{figure}[H]

                            %  \includegraphics[scale=0.5]{images/WhatsApp Image 2021-06-21 at 11.37.55.jpeg}
                            \caption{Table des produits}
                        \end{figure}
                        Chaque produit s'enregistre dans la base des données suivant plusieurs spécifications, l'identifiant du produit, nom , prix ainsi que la présentation visuelle du produit. Avec les informations par rapport à la disponibilité dans le stock et la condition du produit.
                        \subsection{ Table des paniers: }
                        \begin{figure}[H]

                            %  \includegraphics[scale=0.5]{images/WhatsApp Image 2021-06-21 at 11.37.22.jpeg}
                            \caption{Table des paniers}
                        \end{figure}
                        Les paniers fonctionnent comme des paniers réelles, ce tableau fait la spécification de la quantité des produits dans le panier,  à l'aide des identificateurs panier et produits pour bien déterminer le positionnement des produits dans les paniers.
                        \subsection{ Table des achats: }
                        \begin{figure}[H]

                            %  \includegraphics[scale=0.5]{images/WhatsApp Image 2021-06-21 at 11.36.25.jpeg}
                            \caption{Table des achats}
                        \end{figure}
                        En ce qui concerne les ordres d'achats enregistrés, table  des achats fait l'inscription de tout les données concernant les ordres, elle désigne l'adresse mail du client son adresse domicile, la  date dont l'ordre a été effectué en plus du prix totale de l'ordre.
                        \subsection{ Table clients: }
                        \begin{figure}[H]

                            %  \includegraphics[scale=0.5]{images/WhatsApp Image 2021-06-21 at 11.35.41 (1).jpeg}
                            \caption{Table clients}
                        \end{figure}
                        La table client comme son nom l'indique contient tout les indicateurs du client son nom, la date de la création de son compte client ainsi que son Id-utilisateur.

                        \section{Création des pages de l'application web:}
                        \subsection{Page d'acceuil:}
                        \begin{figure}[H]
                            \begin{center}
                                %  \includegraphics[scale=0.7]{images/a1.PNG}
                                \caption{Page d'acceuil}
                            \end{center}
                        \end{figure}
                        \newpage
                        La page d'acceuil expose une selection des produits, avec leurs images réelles, chaque case contient un produit différent avec la possibilité de l'ajouter en panier. En haut de la page on trouve le logo significatif, la barre en haut donne la possibilité d'avoir un panier ou faire le login aux comptes.
                        \begin{figure}[H]
                            \begin{center}
                                %  \includegraphics[scale=0.7]{images/a2.PNG}
                                \caption{Page d'acceuil section aide }
                            \end{center}
                        \end{figure}
                        La figure indique les conseils à suivre pour bien maintenir les produits, ainsi qu'elle offre aux utilisateurs la possibilité de contacter l'administrateur pour signaler leurs problèmes.
                        \subsection{Page authentifiation adminstrateur:}
                        \begin{figure}[H]
                            \begin{center}
                                %  \includegraphics[scale=0.3]{images/WhatsApp Image 2021-06-21 at 17.59.39 (2).jpeg}
                                \caption{Page authentifiation adminstrateur}
                            \end{center}
                        \end{figure}
                        L’authentification est une étape simple, l'administrateur insert les données déjà inscrites dans la base des données, il aura accès à son compte quand ses données sont valables.
                        \subsection{Page authentification adminstrateur:}
                        \begin{figure}[H]
                            \begin{center}
                                %  \includegraphics[scale=0.3]{images/WhatsApp Image 2021-06-21 at 17.59.39 (2).jpeg}
                                \caption{Page authentification adminstrateur}
                            \end{center}
                        \end{figure}
                        \subsection{Page espace adminstrateur:}
                        \begin{figure}[H]
                            \begin{center}
                                %  \includegraphics[scale=0.3]{images/WhatsApp Image 2021-06-21 at 17.59.39 (4).jpeg}
                                \caption{Page espace administrateur}
                            \end{center}
                        \end{figure}
                        C'est l'espace administrateur où il pourra avoir l'historique de toutes les actions, les produits vendus, l'état des paniers et des ordres des clients.
                        \subsection{Graphe Graphe nombre d'utilisateur:}

                        \begin{figure}[H]
                            \begin{center}
                                %  \includegraphics[scale=0.3]{images/WhatsApp Image 2021-06-22 at 00.49.13.jpeg}
                                \caption{Graphe nombre d'utilisateur}
                            \end{center}
                        \end{figure}
                        Dans cette figure on peut savoir à tout moment l'efficience de l'application grace à ce type de graphe qui indique le taux d'utilisateurs qui utilisent l'application par jours .
                        \subsection{Graphe prévision des ventes:}
                        Pour une gestion plus fiable et moins couteuse on a choisit de faire une prévision d'avance qui sert à prévenir les achats qui peuvent être effectués à la période d'un an, en utilisant la méthode de regression linéaire, pour diminuer les pertes qui influencent le capitale et la satisfaction des clients.
                        \begin{figure}[H]
                            \begin{center}
                                %  \includegraphics[scale=0.6]{images/WhatsApp Image 2021-06-22 at 01.31.37.jpeg}
                                \caption{Graphe prévision des ventes}
                            \end{center}
                        \end{figure}
                        \subsection{Page produit type PC Portable:}
                        \begin{figure}[H]
                            \begin{center}
                                %  \includegraphics[scale=0.7]{images/T1.PNG}
                                \caption{Page produit type PC Portable}
                            \end{center}
                        \end{figure}
                        Cette page affiche tout les produits de type ordinateurs portables, en accordant le prix de chaque produit, sa condition et sa fiche informative pour mieux connaitre le produit.
                        \subsection{Page produit type accessoire:}
                        \begin{figure}[H]
                            \begin{center}
                                %  \includegraphics[scale=0.7]{images/T2.PNG}
                                \caption{Page produit type accessoire:}
                            \end{center}
                        \end{figure}

                        La plateforme offre une diversité des accessoires souris, cables de types qui diffèrent selon l'utilité, chaque case donne la description technique des accessoires ainsi que la possibilité de les ajouter au panier pour effectuer l'achat.

                        \subsection{Page produit type stockages:}
                        \begin{figure}[H]
                            \begin{center}
                                %  \includegraphics[scale=0.7]{images/T3.PNG}
                                \caption{Page produit type stockage:}
                            \end{center}
                        \end{figure}

                        Les accessoires du stockage font aussi parties des produits exposés avec leurs deux types SDD et HDD, les accessoires stokages sont réparties dans des cases avec le prix qui leurs correspond et l'image du produit.
                        \subsection{Recapitulatif:}
                        La réalisation du projet est un chapitre très descriptive et abrégeant toutes les étapes réelles de la mise en œuvre de l'application web. En premier, il aborde toutes les outils utilisés tel que le langage  Python. Ensuite, ce chapitre en son intégralité mentionne toutes technologies utisées pour la création des interfaces dynamiques. Finalement une grande partie, consacrée à la consultation de toutes les parties du site et de tous les cas de l’utilisation du « Laptop+ », avec l’insertion des screenshots qui simulent toutes les utilisations possibles avec la totalité des fonctionnalités offertes par l'application ainsi que celles du graphe de prévision qui sert à la bonne gestion des ventes au sein de l'application.
                        \chapter{Conclusion générale et perspectives}
                        \fancyhead[R]{\textbf{Chapitre \thechapter:Conclusion générale et perspectives}}
                        \fancyhead[L]{\hspace*{5cm}}
                        \section{Conclusion générale}
                        L'e-commerce n'est pas la thématique qu'on aura apprendre seulement sous un contexte académique, mais un état d'esprit de mise en question et de prise de decision dans un cadre entreprenariale qui soutient la digitalisation de plusieurs secteurs. Cela explique notre choix de sujet su lequel on travaillait en particulier ces derniers mois, sous la demande énorme de l'informatisation des opérations d'achat effectuées par les consommateurs dans l'agenda des réglementations conçues sous le cadre de la pandémie.

                        Actuellement, les documentations montrent que les achats en ligne sont plus utiles de la part des clients et pour les firmes en terme des revenus importants, ce qui les distingue comme les points que  l'application vise en premier.

                        Notre idée part d'un contexte digitale qui s'intitule \textbf{Laptop+}: C’est une application à utiliser par les entreprises spécialisées par la vente des ordinateurs, accessoires et des outils de stockage, pour encourager les client à utiliser ce mode d'achat, en leurs offrant une solution pratique, fiable et en plus  sécurisé en terme des transactions et des paiements.

                        Ce projet a été une belle occasion pour travailler en groupe en des conditions qui ne sont pas assez simples, pour développer nos
                        connaissances et nos compétences et surtout pour découvrir de nouvelles fonctionnalités des plateformes d'achats en ligne et des prévision des ventes et des langages tels que Python , la création des graphe dans la phase technique.
                        \section{Perspectives}
                        \chapter{Webographie}
                        \fancyhead[R]{\textbf{Chapitre \thechapter:Webographie}}
                        \fancyhead[L]{\hspace*{5cm}}
                        \section{Plateforme open source:}
                        \begin{itemize}
                            \item \textbf{[kaggle]:} https://www.kaggle.com/carrie1/ecommerce-data, dernière mise à jour le 15/06/2021.
                            \item \textbf{[GitHub]:} https://github.com/topics/ecommerce-website, dernière mise à jour le 20/06/2021.
                        \end{itemize}
                        \section{Cours:}
                        \begin{itemize}
                            \item \textbf{Coursera:} https://www.coursera.org/specializations/django, dernière mise à jour le 20/06/2021 mais avec inscription.
                            \item\textbf{Udemy:}https://www.udemy.com/course/learn-python-django-from-scratch/, dernière mise à jour 21/06/2021.
                            \item \textbf{Citycours :}https://sites.uclouvain.be/P2SINF/chartjs.html, dernière mise à jour le 21/06/2021.
                        \end{itemize}
                        \section{Dataset:}
                        \begin{itemize}
                            \item \textbf{Kaggle:} https://www.kaggle.com/jzeferino/predictive-sales-with-linear-regressions, dernière mise à jour le\\ 19/06/2021.
                        \end{itemize}
                    \end{doublespace}
                \end{doublespace}
            \end{doublespace}
        \end{doublespace}
    \end{doublespace}
\end{doublespace}
\end{document}